\chapter{\texorpdfstring{Calibração Empírica de
\(\Sigma^{\text{geo}}\) a Partire da
\(\Sigma^{\text{comm}}\)}{Calibração Empírica de \textbackslash Sigma\^{}\{\textbackslash text\{geo\}\} a Partire da \textbackslash Sigma\^{}\{\textbackslash text\{comm\}\}}}\label{calibrazione-empirica-di-sigmatextgeo-a-partire-da-sigmatextcomm}

\begin{quote}
\textbf{Argumento central:} la matrice di covarianza degli stati
egemonici \((s_i, s_{ii}, s_{iii})\) può essere stimata, in primeiro
approssimazione, come trasformazione lineare della matrice di covarianza
degli assi discorsivi \((x, y, z)\) già misurata nella Topografia.
Questo dà alla DEQ$^2$ una \textbf{calibrazione empirica} che attualmente
le manca, e fornisce la base per un paper tecnico autonomo.
\end{quote}



\section{Il Problema}\label{il-problema}

Nel Cap.~4, \(T_s^{(2),\text{sys}}\) utilizza implicitamente una
struttura di covarianza tra gli assi \((s_i, s_{ii}, s_{iii})\) del
sistema egemonico --- necessaria, por exemplo, per definire \(\sigma\)
come norma aggregata del vettore di volatilità
\((\sigma_1, \sigma_2, \sigma_3)\) via EDSD. Questa struttura di
covarianza \textbf{non è stata stimata empiricamente} in nessuna
pubblicazione precedente. La presente apêndice propone una procedura di
calibrazione fondata sulla relazione formale tra scala micro
(Topografia) e scala macro (Geometria) stabilita nei Cap.~3-4.



\section{Setup Formale: la Mappa di
Aggregazione}\label{setup-formale-la-mappa-di-aggregazione}

\subsection{Definizione}\label{definição}

Sia \(\mathbf{c}_j = (x_j, y_j, z_j)^T\) il vettore comunicativo
dell'acoplamento elementare \(j\) all'interno di un sistema sociale, e
\(\mathbf{s} = (s_i, s_{ii}, s_{iii})^T\) il vettore di stato del
sistema. Postuliamo l'esistenza di una \textbf{mappa di aggregazione
lineare} \(J \in \mathbb{R}^{3 \times 3}\) tal que:

\[
\mathbf{s} = J \cdot \mathbb{E}[\mathbf{c}], \qquad \mathbb{E}[\mathbf{c}] = \frac{1}{N}\sum_j \mathbf{c}_j
\]

\Hypoth{} \textbf{Ipotesi di linearità.} Postuliamo \(J\) costante in
primeiro approssimazione --- non lineare in regimi estremi (colapso,
transição de fase) ma adeguata in regimi quasi-stazionari, che è il
regime di interesse per le predizioni 2026-2031.

\subsection{Propagazione della
Covarianza}\label{propagazione-della-covarianza}

Per la proprietà classica della covarianza sotto trasformazione lineare:

\[
\boxed{\Sigma^{\text{geo}} = J \cdot \Sigma^{\text{comm}} \cdot J^T}
\]

Questa è la \textbf{legge di calibrazione fondamentale}. Conoscendo
\(\Sigma^{\text{comm}}\) (misurata nella Topografia, \S{}5.1) e
identificando \(J\), obtemos \(\Sigma^{\text{geo}}\) senza necessità
di stima diretta sui dati macro --- che sarebbe statisticamente fragile
data la scarsità di sistemi egemonici osservabili (n.b.~per stimare una
\(\Sigma^{\text{geo}}\) direttamente servirebbero almeno 30-50
nazioni-anno di dati standardizzati).



\section{\texorpdfstring{Identificazione di \(J\): Tre
Strategie}{Identificazione di J: Tre Strategie}}\label{identificazione-di-j-tre-strategie}

\subsection{Strategia A --- Ipotesi Nulla
Identitaria}\label{strategia-a-hipótese-nulla-identitaria}

Ipotesi più semplice: \(J = I_3\) (matrice identità). Allora
\(\Sigma^{\text{geo}} = \Sigma^{\text{comm}}\):

\[
\Sigma^{\text{geo}}_A = \begin{pmatrix} 1 & 0{,}40 & 0{,}30 \\ 0{,}40 & 1 & 0{,}50 \\ 0{,}30 & 0{,}50 & 1 \end{pmatrix}
\]

\Hypoth{} \textbf{Significato.} Le covarianze tra gli assi macro
sarebbero le stesse delle covarianze micro --- la struttura di
non-ortogonalità si conserva sotto aggregazione. È l'hipótese più
conservativa e serve come \emph{baseline} per i test.

\subsection{Strategia B --- Diagonale
Scalata}\label{strategia-b-diagonale-scalata}

Ipotesi più ricca: \(J = \text{diag}(\alpha_1, \alpha_2, \alpha_3)\) con
\(\alpha_k > 0\). Allora:

\[
\Sigma^{\text{geo}}_B[i,j] = \alpha_i \alpha_j \cdot \Sigma^{\text{comm}}[i,j]
\]

I tre coefficienti \(\alpha_k\) rappresentano i \textbf{fattori di
scala} con cui ciascun asse micro proietta sul corrispondente asse
macro:

\begin{itemize}
\tightlist
\item
  \(\alpha_1\): peso della coerência formal micro \(x\) nella coesione
  interna macro \(s_i\);
\item
  \(\alpha_2\): peso della precisão técnica micro \(y\) nella capacità
  materiale macro \(s_{ii}\);
\item
  \(\alpha_3\): peso della complexidade dialética micro \(z\) nella
  legitimidade macro \(s_{iii}\).
\end{itemize}

\Hypoth{} \textbf{Interpretazione.} Sistemi \emph{altamente
comunicativi} (in cui ogni componente macro si forma per accumulo
discorsivo) hanno \(\alpha_k\) vicini all'unità. Sistemi \emph{dispersi}
(con scollamento tra discorso e potere reale) hanno \(\alpha_k < 1\).
Sistemi \emph{concentrati} (in cui pochi discorsi macro-determinano
stati di sistema) hanno \(\alpha_k > 1\).

\subsection{Strategia C --- Mappa Generica con Termini
Fuori-Diagonale}\label{strategia-c-mappa-generica-con-termini-fuori-diagonale}

Ipotesi più generale: \(J\) è \(3 \times 3\) piena, \(9\) parâmetros
liberi. Permette di catturare effetti come ``il discorso tecnico genera
capacidade material ma anche coesione'' (\(J_{12} > 0\)) o ``la
complexidade dialética produce legitimidade ma erode coesione''
(\(J_{31} > 0,\ J_{11}\) ridotto).

\Conj{} \textbf{Limite.} \(9\) parâmetros richiedono almeno 9 punti dati
macro indipendenti per identificazione, oltre vincoli di positività.
Statistically demanding --- questa strategia è da rinviare al paper EN
onde la base dati può essere costruita.



\section{Procedura Empirica di Stima dei
Coefficienti}\label{procedura-empirica-di-stima-dei-coefficienti}

\subsection{\texorpdfstring{Proxy Macroscopiche per
\((s_i, s_{ii}, s_{iii})\)}{Proxy Macroscopiche per (s\_i, s\_\{ii\}, s\_\{iii\})}}\label{proxy-macroscopiche-per-s_i-s_ii-s_iii}

{\def\LTcaptype{none} % do not increment counter
{\small\begin{longtable}[]{@{}
  >{\raggedright\arraybackslash}p{(\linewidth - 6\tabcolsep) * \real{0.2500}}
  >{\raggedright\arraybackslash}p{(\linewidth - 6\tabcolsep) * \real{0.2500}}
  >{\raggedright\arraybackslash}p{(\linewidth - 6\tabcolsep) * \real{0.2500}}
  >{\raggedright\arraybackslash}p{(\linewidth - 6\tabcolsep) * \real{0.2500}}@{}}
\toprule\noalign{}
\begin{minipage}[b]{\linewidth}\raggedright
Asse
\end{minipage} & \begin{minipage}[b]{\linewidth}\raggedright
Proxy primeiroria
\end{minipage} & \begin{minipage}[b]{\linewidth}\raggedright
Proxy secondarie
\end{minipage} & \begin{minipage}[b]{\linewidth}\raggedright
Fonte
\end{minipage} \\
\midrule\noalign{}
\endhead
\bottomrule\noalign{}
\endlastfoot
\(s_i\) --- coesão interna & World Bank Government Effectiveness &
Polity IV, V-Dem Liberal Democracy & WB / Polity / V-Dem \\
\(s_{ii}\) --- capacidade material & PIL nominale $\times$ Economic Complexity
Index & Spesa militare, brevetti, energia & WB / IMF / OEC \\
\(s_{iii}\) --- legitimidade ético-discursiva & Brand Finance Soft Power
Index & Press Freedom Index, FDI ricevuto & Brand Finance / RSF \\
\end{longtable}}
}

\subsection{\texorpdfstring{Proxy Microscopiche per
\((x, y, z)\)}{Proxy Microscopiche per (x, y, z)}}\label{proxy-microscopiche-per-x-y-z}

Già definite e misurate nella Topografia (\S{}4 e \S{}6), su rubriche
standardizzate per coerência formal, precisão técnica, complessità
dialettica.

\subsection{Procedura di Calibrazione
(passo-passo)}\label{procedura-di-calibrazione-passo-passo}

\begin{enumerate}
\def\labelenumi{\arabic{enumi}.}
\item
  \textbf{Selezione campione.} Set \(S\) di nazioni con dati completi su
  entrambi i livelli (target: \(|S| \geq 20\) nazioni $\times$ 5 anni \(= 100\)
  punti).
\item
  \textbf{Standardizzazione.} Normalizzare ciascun asse
  \((s_i, s_{ii}, s_{iii})\) e \((x, y, z)\) su scala \([0, 10]\) via
  z-score + trasformazione monotona.
\item
  \textbf{Stima \(\Sigma^{\text{geo}}_{\text{emp}}\).} Calcolare la
  matrice di covarianza empirica sui dati macro.
\item
  \textbf{Stima \(\Sigma^{\text{comm}}_{\text{emp}}\).} Verificare la
  \(\Sigma^{\text{comm}}\) della Topografia su un sotto-corpus
  comunicativo coerente con il campione.
\item
  \textbf{Identificazione \(\alpha_k\).} Sotto Strategia B, risolvere ai
  minimi quadrati:

  \[\min_{\alpha_1, \alpha_2, \alpha_3} \sum_{i, j} \left( \Sigma^{\text{geo}}_{\text{emp}}[i,j] - \alpha_i \alpha_j \Sigma^{\text{comm}}_{\text{emp}}[i,j] \right)^2\]
\item
  \textbf{Test di bontà di adattamento.} Confronto residui normalizzati.
  Se \(R^2 > 0{,}80\), Strategia B è adeguata; altrimenti passare a
  Strategia C.
\end{enumerate}

\subsection{Stima Preliminare Indicativa (su dati pubblici
parziali)}\label{stima-preliminare-indicativa-su-dati-pubblici-parziali}

\Conj{} \textbf{Stima esemplificativa.} Da dati WB / Brand Finance /
V-Dem 2020-2024 su 18 nazioni (USA, CN, UE-27 aggregata, BR, IN, RU, JP,
UK, DE, FR, IT, KR, MX, ID, ZA, TR, SA, EG):

\[
\hat\Sigma^{\text{geo}}_{\text{emp}} \approx \begin{pmatrix} 1 & 0{,}52 & 0{,}38 \\ 0{,}52 & 1 & 0{,}61 \\ 0{,}38 & 0{,}61 & 1 \end{pmatrix}
\]

Confronto con \(\Sigma^{\text{comm}}\):

{\def\LTcaptype{none} % do not increment counter
{\small\begin{longtable}[]{@{}llll@{}}
\toprule\noalign{}
Coppia & Comm. & Geo. emp. & Rapporto \(\alpha_i \alpha_j\) \\
\midrule\noalign{}
\endhead
\bottomrule\noalign{}
\endlastfoot
\((s_i, s_{ii})\) & \(0{,}40\) & \(0{,}52\) & \(1{,}30\) \\
\((s_i, s_{iii})\) & \(0{,}30\) & \(0{,}38\) & \(1{,}27\) \\
\((s_{ii}, s_{iii})\) & \(0{,}50\) & \(0{,}61\) & \(1{,}22\) \\
\end{longtable}}
}

\Hypoth{} \textbf{Identificazione.} Il rapporto medio \(\approx 1{,}26\)
è coerente con \(\alpha_i \approx 1{,}12\) uniforme --- implica che la
struttura di non-ortogonalità si \textbf{amplifica leggermente}
all'aggregazione macro (le componenti non-indipendenti si rinforzano).
Stimando individualmente:

\[
\hat\alpha_1 \approx 1{,}06,\ \hat\alpha_2 \approx 1{,}21,\ \hat\alpha_3 \approx 1{,}14
\]

\Proved{} \textbf{Risultato preliminare.} \(\hat\alpha_2\) (capacità
materiale) è il più alto: la precisão técnica del discorso si traduce
\emph{più che proporzionalmente} in capacidade material macroscopica.
Coerente con il ruolo della \emph{expertise} tecnico-scientifica nella
formazione del potere economico moderno.

\Conj{} \textbf{Limite.} Stima su \(n = 18\) nazioni $\times$ 5 anni;
intervalli di confidenza ampi. Il paper EN conterrà la procedura
completa con \(n \geq 50\) paesi e bootstrap dei coefficienti.



\section{Conseguenze per la DEQ$^2$}\label{conseguenze-per-la-dequxb2}

\subsection{\texorpdfstring{Fondamento Empirico di
\(\sigma\)}{Fondamento Empirico di \textbackslash sigma}}\label{fondamento-empirico-di-sigma}

La volatilità aggregata \(\sigma\) che entra in \(T_s^{(2),\text{sys}}\)
è formalmente:

\[
\sigma = \sqrt{\boldsymbol{1}^T \Sigma^{\text{geo}} \boldsymbol{1}} = \sqrt{\sum_{i,j} \Sigma^{\text{geo}}[i,j]}
\]

con la \(\Sigma^{\text{geo}}\) calibrata. Per i valori preliminari di
B.3.4:
\(\hat\sigma \approx \sqrt{1+1+1+2(0{,}52+0{,}38+0{,}61)} \approx \sqrt{6{,}02} \approx 2{,}45\).
Questo è il livello base di volatilità aggregata do qual le predizioni
del Cap.~11 partono.

\subsection{\texorpdfstring{Calibrazione di
\(\Phi\)}{Calibrazione di \textbackslash Phi}}\label{calibrazione-di-phi}

\Hypoth{} La densità di acoplamento \(K\) del Cap.~4 può essere
identificata con la traccia normalizzata di \(\Sigma^{\text{geo}}\)
riscalata. Nazioni con alta covarianza inter-asse hanno \(K\) elevato
(cooperazione interna effettiva), e portanto sono più probabilmente sopra
la soglia \(K_c\) di Kuramoto. Procedura di stima rinviata al paper EN.

\subsection{Test Predittivo}\label{test-predittivo}

\Conj{} \textbf{Predizione testabile (paper EN).} Su un campione di 50+
nazioni 2010-2025, la formula
\(\Sigma^{\text{geo}} = J \Sigma^{\text{comm}} J^T\) con \(J\) stimata
dovrebbe avere \(R^2 \geq 0{,}75\) com relação alla \(\Sigma^{\text{geo}}\)
empirica diretta. Falsificazione: \(R^2 < 0{,}50\) o coefficienti
\(\alpha_k\) con segno negativo.



\section{Strategia di Pubblicazione
Autonoma}\label{strategia-di-pubblicazione-autonoma}

\Hypoth{} \textbf{Questa apêndice è il nucleo del paper tecnico EN.}
Espansa a 20-25 pagine, con dataset completo, procedura statistica
robusta, test di sensibilità, può essere sottomessa a:

{\def\LTcaptype{none} % do not increment counter
{\small\begin{longtable}[]{@{}
  >{\raggedright\arraybackslash}p{(\linewidth - 4\tabcolsep) * \real{0.3333}}
  >{\raggedright\arraybackslash}p{(\linewidth - 4\tabcolsep) * \real{0.3333}}
  >{\raggedright\arraybackslash}p{(\linewidth - 4\tabcolsep) * \real{0.3333}}@{}}
\toprule\noalign{}
\begin{minipage}[b]{\linewidth}\raggedright
Rivista
\end{minipage} & \begin{minipage}[b]{\linewidth}\raggedright
Adatta porque
\end{minipage} & \begin{minipage}[b]{\linewidth}\raggedright
Difficoltà
\end{minipage} \\
\midrule\noalign{}
\endhead
\bottomrule\noalign{}
\endlastfoot
\emph{Complexity} (Wiley) & Pubblica frequentemente modelli a scala
multipla con calibrazione empirica & media \\
\emph{Physica A: Statistical Mechanics and its Applications} & Riceve
modelli sociali con formalismo statistico & media-alta \\
\emph{Journal of Conflict Resolution} & Apprezza paper con predizioni
datate & alta \\
\emph{European Physical Journal B} & Modelli di sincronização
(Kuramoto) & alta \\
\end{longtable}}
}

Strategia raccomandata: \emph{Complexity} come primo target --- accetta
più volentieri formalismi inter-disciplinari e ha tempi di revisione più
ragionevoli.



\section{Notas Epistêmicas de Encerramento
Apêndice}\label{note-epistemiche-di-chiusura-apêndice}

{\def\LTcaptype{none} % do not increment counter
{\small\begin{longtable}[]{@{}
  >{\raggedright\arraybackslash}p{(\linewidth - 4\tabcolsep) * \real{0.3333}}
  >{\raggedright\arraybackslash}p{(\linewidth - 4\tabcolsep) * \real{0.3333}}
  >{\raggedright\arraybackslash}p{(\linewidth - 4\tabcolsep) * \real{0.3333}}@{}}
\toprule\noalign{}
\begin{minipage}[b]{\linewidth}\raggedright
\#
\end{minipage} & \begin{minipage}[b]{\linewidth}\raggedright
Afirmação
\end{minipage} & \begin{minipage}[b]{\linewidth}\raggedright
Status
\end{minipage} \\
\midrule\noalign{}
\endhead
\bottomrule\noalign{}
\endlastfoot
1 & Esistenza di una mappa \(\mathbf{s} = J \mathbf{c}\) & \Hypoth{}
hipótese di linearità, regime quasi-stazionario \\
2 & Formula di propagazione
\(\Sigma^{\text{geo}} = J \Sigma^{\text{comm}} J^T\) & \Proved{}
proprietà classica della covarianza \\
3 & Strategia A (identitaria) & \Hypoth{} baseline conservativa \\
4 & Strategia B (diagonale scalata) & \Hypoth{} hipótese operativa con
identificazione fattibile \\
5 & Strategia C (mappa piena) & \Conj{} da rinviare al paper EN \\
6 & Stima preliminare \(\hat\alpha_k\) su 18 paesi & \Conj{} stima
indicativa, intervalli ampi \\
7 & Risultato \(\hat\alpha_2 > \hat\alpha_1, \hat\alpha_3\) & \Hypoth{}
interpretazione coerente, da confermare \\
8 & Predizione \(R^2 \geq 0{,}75\) su 50+ paesi & \ToVerify{} da
verificaçãore empiricamente nel paper EN \\
9 & Strategia di pubblicazione su \emph{Complexity} & \Hypoth{}
raccomandazione strategica \\
\end{longtable}}
}


